%%=============================================================================
%% 中英文目录
%%=============================================================================

%% 中英文目录
\addtocontents{toc}{\protect\hypersetup{hidelinks}}
\addtocontents{etoc}{\protect\hypersetup{hidelinks}}

%% 中文目录
\renewcommand\tableofcontents{
	\setstretch{1} % 设置行间距

	\chapter*{\contentsname}

	\ifxmu@compact
		\markboth{\xmu@label@compactcontents}{\xmu@label@compactcontents} % 页眉
		\addcontentsline{toc}{chapter}{\xmu@label@compactcontents} %将中文摘要加入到目录中
	\else
		\markboth{\xmu@label@contents}{\xmu@label@contents} % 页眉
		\addcontentsline{toc}{chapter}{\xmu@label@contents} %将中文摘要加入到目录中
	\fi

	% 	\addtocontents{etoc}{\protect\contentsline{chapter}{\sffamily \encontentsname}{\sffamily \thepage}{chapter.\thechapter}} %将 Contents 加入到英文目录中

	\addcontentsline{etoc}{chapter}{\sffamily \xmu@label@encontentszh} %将中文目录加入到英文目录中
	% 	\addtocontents{etoc}{\protect\contentsline{chapter}{\sffamily \encontentsname}{\thepage}{chapter.\thechapter}} %将 Contents 加入到英文目录中

	\makeatletter
	\renewcommand{\@pnumwidth}{0em}
	\makeatother
	\@starttoc{toc}%
	\clearpage
}
\setcounter{secnumdepth}{4}  % 目录显示的章节编号深度 (part 对应 -1)
\setcounter{tocdepth}{2}     % 目录深度 (part 对应 -1)

%% 添加英文目录条目的命令
\newcommand\enchapter[1]{\addtocontents{etoc}
	{\protect\contentsline{chapter}
		{\protect\numberline{\sffamily Chapter \thechapter}#1}
		% 		{\sffamily \thepage}{chapter.\thechapter}}
		{\thepage}{chapter.\thechapter}}
}
\newcommand\ensection[1]{\addtocontents{etoc}
	{\protect\contentsline{section}
		{\protect\numberline{\sffamily \thesection}#1}
		% 		{\sffamily \thepage}{section.\thesection}}
		{\thepage}{section.\thesection}}
}
\newcommand\ensubsection[1]{\addcontentsline{etoc}{subsection}{\protect\numberline{\thesubsection}#1}}
\newcommand\ensubsubsection[1]{\addcontentsline{etoc}{subsubsection}{\protect\numberline{\thesubsubsection}#1}}

%% 双语文档结构宏（同时写入中文和英文目录）
\newcommand\xmuchapter[2]{\chapter{#1}\enchapter{#2}}
\newcommand\xmusection[2]{\section{#1}\ensection{#2}}
\newcommand\xmusubsection[2]{\subsection{#1}\ensubsection{#2}}

%% 英文目录
\newcommand\entableofcontents{
	\setstretch{1} % 设置行间距

	\chapter*{\sffamily\encontentsname}

	\pdfbookmark[0]{\encontentsname}{english_contents}

	\markboth{\encontentsname}{\encontentsname} % 修改页眉文字

	\addcontentsline{toc}{chapter}{\xmu@label@encontents} % 将英文目录加入到中文目录中
	\addcontentsline{etoc}{chapter}{\sffamily \xmu@label@encontentsen} % 将英文目录加入到英文目录中

	\makeatletter
	\renewcommand{\@pnumwidth}{0em}
	\makeatother

	\@starttoc{etoc}%
	\clearpage
}
